% Cluster F: scope and limitations (Warning 10.59 discipline)
% Body-only LaTeX; preamble macros, theorem environments, and cite keys per CONVENTIONS.md.

\section{Scope and limitations}\label{sec:scope}

The results of this paper form a closed package, and the manuscript
\cite{MWman} disciplines its use with an explicit warning (its
Warning~10.59), whose discipline we retain here in full.  The package has
exactly five load-bearing statements, and their division of labor should
be read precisely.  Theorem~\ref{thm:realization} proves realization
\emph{from} a certificate: given a protected ribbon-routing certificate,
it produces a compact connected full carrier with singleton attachment
traces, surviving protected chains, injective survival, and preserved
cyclic order.  Theorem~\ref{thm:free-component} characterizes \emph{when}
such a certificate exists: exactly when one component of the frozen free
domain has simultaneous singleton ports to every required atom, with the
persistence clause that a separation visible in the frozen data can never
be repaired by monotone connector enlargement.
Theorem~\ref{thm:common-domain} is \emph{one sufficient planar route} to
that criterion: an augmented-atom packet system forces a single common
free component.  Theorem~\ref{thm:audit} removes completion order and
global port untangling from that route: the packet completions are
order-independent, the port checks localize, and merger, localization,
and port-loss defects are persistent certificates of failure.  Finally,
Theorem~\ref{thm:carrier} instantiates the entire package at finite
direct-decoration scale, where the cells of the atlas play the role of
the cellwise data of Proposition~\ref{prop:cellwise} and the nested
packet lemma absorbs the infinitely many unselected fibers.

The mathematical ingredients of all five statements are classical:
Janiszewski's separation theorem and the plane topology of continua
\cite{New51,Why42}, prime ends \cite{Pom92}, filled hulls, and
\v{C}ech cohomology with the compacta form of Alexander duality
\cite{Hat02}, in the tradition of planar arguments long standard in
complex dynamics \cite{DH85}.  The contribution is not a new
theorem about any of these ingredients; it is the assembled
certification framework---explicitly protected accesses with declared
sides, a frozen future inventory fixed before any connector is chosen,
typed packets with order-independent completions, and monotone,
persistent failure certificates.  The paper is proof technology, and its
value is exactly the conversion it performs: informal planar-routing
arguments of the form ``connect these continua while keeping these
accesses and avoiding that set,'' which are traditionally re-derived ad
hoc inside larger proofs, become exact, persistent, checkable
certificates that can be audited clause by clause, transported by
univalent maps, and reused without re-proof.

\begin{warning}\label{warn:F-scope}
Theorem~\ref{thm:realization} proves realization from a certificate.
Theorem~\ref{thm:free-component} characterizes when that certificate
exists, Theorem~\ref{thm:common-domain} proves one sufficient planar
route, and Theorem~\ref{thm:audit} removes completion order and port
untangling from that route.  Theorem~\ref{thm:carrier} instantiates the
package at finite direct-decoration scale, and it is finite and
direct-scale only (Warning~\ref{warn:F-finite-scale}).  None of these
results constructs the virtual Cantor package, the complete Wave
horizon, the analytic covering or immersion maps, a period-independent
incidence bound or future-faithful compression, or any width,
compactness, rigidity, or MLC conclusion.  The
open problems of the source manuscript---10.6, 10.50--10.52, 10.55,
10.57, and 10.58 in the numbering of \cite{MWman}---remain open and are
not advanced here: the virtual canonical interface carrier, the
manuscript's first protected-tower clause, and MLC itself remain open.
\end{warning}

For the reader consulting the manuscript, the open problems named in the
warning are, briefly: 10.6, the virtual protected routing certificate for
the lamination-lift candidate carrier; 10.50, the virtual cylinder
target-address selector; 10.51, the transfer of the disconnected marking
to actual bases and actual iterates on the cylinders and end fibers of
the pullback compactification; 10.52, the virtual Cantor-package
realization at every intermediate scale; and 10.58, the virtual common
free component through the complete Wave horizon, together with the
related problems 10.55 and 10.57.  On the last of these the manuscript is explicit
that neither the connected-hull construction of \cite{Dud25} nor the
primitive-category Hubbard-tree results of \cite{KKL26} supply the
completed free domain that the problem demands; nothing in the present
paper does so either, and nothing here is conditional on those problems:
every theorem above is proved outright in its finite planar setting.

Two further boundaries deserve restatement.  First, the only intentional
divergence from the printed manuscript is the repair of Lemma~10.7 and
the correspondingly rewritten proof of Theorem~10.8, recorded in
Remark~\ref{rem:repair}; every other statement is a faithful
transcription.  Second, no result of this paper asserts local
connectivity of any atom, of any fiber, or of any Julia set, and no
result should be quoted as progress toward MLC.  What may be quoted is
the technology: certificates that realize, criteria that are exact,
audits that are order-free, and one finite dynamical instantiation in
which all of it is checkable.
